\vspace{-0.5em}
\section{Graph Query Processing}
\label{sec:graph_qp}

Having described the system architecture and the storage model, we now
discuss challenges introduced by executing typical graph database queries 
obliviously.

\vspace{-1em}
\subsection{0-hop query}
\label{sub:0-hop}

We begin by examining the simplest type of query in graph databases involving a single 
node label with one or more predicates on its properties. This section 
highlights the challenges posed by non-oblivious query processing algorithms and provides insights 
on how to morph these algorithms to ensure obliviousness. To illustrate, let's take a sample Cypher query:
\texttt{MATCH u:user WHERE u.age $>$ 30 AND u.lives\_in = `LA' RETURN u.name}.

Consider the following code that retrieves the properties relevant to the query
and tests the predicates using an if-else statement.

\noindent\hrulefill
\begin{algorithmic}[1]
        \State Assume $ages$, $lives\_in$, and $names$ are lists that store relevant properties after reading from the storage
        \State $output \gets []$
        \For {$i$ in range(0, $num\_users$)}
            \If {$ages[i] > 30$ \&\& $lives\_in[i] == ``LA"$ } 
                \State $output.append(names[i])$
            \EndIf
        \EndFor
        \State \Return $output$
\end{algorithmic}
\noindent\hrulefill

\noindent Executing the above code within an enclave reveals the node ids of users that
satisfy the predicates to an adversary who can observe memory accesses because the 
$output$ variable is accessed and modified only when the predicates are true. 
One way to execute the above functionality obliviously is to use a ternary
operator:

\noindent\hrulefill
\begin{algorithmic}[1]
        \State $output \gets []$
        \For {$i$ in range(0, $num\_users$)}
            \State \parbox[t]{\linewidth}{$res \gets (ages[i] > 30$ \&\& $lives\_in[i] == ``LA"$)?\newline \null\hspace{9mm} $names[i]$ : $``dummy"$} 
            \State $output.append(res)$
        \EndFor
        \State \Return obliviousCompact($output$)
\end{algorithmic}
\noindent\hrulefill

This code modifies the $output$ for each user, ensuring identical 
memory accesses for all node ids. The $output$ can then be 
compacted obliviously using an \textit{oblivious compaction} 
mechanism, such as~\cite{goodrich2011data, sasy2022fast}, to remove dummies before 
returning the result to the client.
Existing oblivious 
databases~\cite{eskandarian2019oblidb, zheng2017opaque, mavrogiannakis2025obliviator, krastnikov2020efficient} employ 
such strategies for processing relational queries. 

\subsection{1-hop query}

Having used the simple 0-hop query to provide intuitions on processing queries
obliviously, this section discuss challenges and future research directions
stemming from 1-hop queries. A 1-hop query spans two distinct node labels $n1$ and $n2$ (we discuss cyclic queries in~\cref{sub:cyclic})
connected by an edge $e$ of the form: $n1 - e \rightarrow n2$, with predicates on any node
or edge label.

\noindent\textit{\textbf{1. Challenges with processing 1-hop query using
existing relational operators.}} One way to process such queries is to consider the two nodes and the edge as relations and use 
oblivious binary join operators~\cite{mavrogiannakis2025obliviator,krastnikov2020efficient, eskandarian2019oblidb,zheng2017opaque}: first join $n1$ with $e$, then join the
output of that with $n2$. Because the edge relation exhibits primary-key-foreign-key relationship with both the node relations, the output of both the joins will
be $|e|$, the number of edges. 

Naively executing a sort-merge join, the most adopted join technique in oblivious settings
~\cite{zheng2017opaque, eskandarian2019oblidb, krastnikov2020efficient, mavrogiannakis2025obliviator},
reveals how many times each join key is repeated based on the pointer traversal pattern 
in the merge step. Existing solutions counter this leakage by first concatenating the two 
input tables and obliviously sorting it on the $<$join key, table id$>$ pair, which in our case, sorts a table of $|n_i| + |e|$ size. This sorted concatenated table
helps identify the join keys present in both tables and how many times each key appears in
the output based on the product of its count in each table.

Applying any of these join algorithms~\cite{zheng2017opaque, eskandarian2019oblidb, krastnikov2020efficient, mavrogiannakis2025obliviator} to process a 1-hop query requires performing two joins, each involving an oblivious sort over tables of sizes 
$|n1|+|e|$ and $|e|+|n2|$. Since oblivious sorting dominates the cost of oblivious joins with 
a complexity of $O(nlog^2n)$ for a table of size $n$, this approach incurs significant 
overhead when $|n|$ and $|e|$ scale to millions of nodes and billions of edges.
% Recall from \cref{sub:storage-layer} that graph databases store the node and edge information already in sorted
% format for better query efficiency.
% But existing oblivious join 
% operators do not take advantage of this because \textit{they all require merging the
% input tables and obliviously sorting it on the join key to identify join keys
% present in both tables}.
This raises the first research question:

\begin{tcolorbox}[colback=gray!7!white,coltitle=black,colframe=orange!50!white,title=\centering \textbf{Research question 1}]
\centering
\emph{Can we design a new oblivious binary join operator that identifies join
keys present in both tables \textbf{without} concatenating and sorting the two input tables?}
\end{tcolorbox}

\noindent\textit{\textbf{Towards a solution}}: One way to identify join keys
present in both tables is to hash the join keys of one table $T1$ and store it 
in a hash map and probe the second table $T2$ by hashing the join keys and checking if
the hashed value is present in the hash table (i.e., a hash join). This solution 
identifies join keys
present in both tables without requiring to concatenate and sort them. However, skew
in the join key distribution reveals sensitive information to an adversary:
if $T1$ has many repeated join keys, this will create hash collisions, revealing
the skew in the data. Alternatively, repeated join keys in $T2$ will repeatedly
probe the same hash bucket, also revealing the skew in the
data. Thus, the oblivious join scheme should protect against such leakages by
considering potential solutions such as key de-duplication and oblivious aggregation operations. 
% Such a join operator will be applicable to many applications where
% the two inputs to the operator come pre-sorted.

\vspace{0.7em}
\noindent\textit{\textbf{2. Challenges with relying on relational operators.}} 
Even with a new efficient oblivious binary join protocol that removes the need
to sort a large concatenated table, executing a 1-hop query using relational
operators incurs high latency because the two joins must be executed \textit{sequentially}.
% Additionally, the most efficient oblivious binary join~\cite{mavrogiannakis2025obliviator},
% which is a sort-merge join algorithm, requires an intermediary sort of the size of the 
% output, which in our case is
% $|e|$. This sort is necessary to align the join keys of the two tables before the 
% final merge phase~\cite{krastnikov2020efficient,mavrogiannakis2025obliviator}.
% Hence, with two joins, this approach sorts two intermediary tables of size $|e|$,
% \textit{sequentially}.
This raises the next research question:

\begin{tcolorbox}[colback=gray!7!white,coltitle=black,colframe=orange!50!white,title=\centering \textbf{Research question 2}]
\centering
\emph{Can we design a custom \textbf{parallelizable} oblivious query processing algorithm to serve 1-hop queries without relying on oblivious binary joins?}
\end{tcolorbox}

\noindent\textit{\textbf{Towards a solution}}: In designing a new oblivious algorithm
to avoid the above sequential processing, we look at plaintext graph databases
that process such queries from \textit{two directions} in parallel: $n1\rightarrow e$
and $e \rightarrow n2$. At a high level, find all nodes ids
of $n1$ that act as the source of an edge with id $e_i$, and node ids of $n2$ that
is the destination of an edge $e_j$. This gives us two lists of edge ids; the 
intersection of these ids is the output of the query (note that we can propagate node
or edge properties along the way). However, naively executing this approach may leak 
sensitive information such as
% \rjm{I think you want to say "the number of nodes with an edge of type e", right?  Or do you me the list of nodes themselves are leaked?}  
the number of unique nodes from $n1$ and $n2$ with an edge label $e$ and the edges in the final output. Hence, this creates the need for custom 1-hop operator 
that traverses the forward and reverse edge lists in parallel to compute the output 
without revealing any sensitive information.

\vspace{-0.7em}
\subsection{2-hop queries and beyond}

Next, we discuss the challenges introduced by multi-hop queries
spanning three or more nodes, connected by two or more
edges.

\noindent\textit{\textbf{Challenges with intermediary result size leakages}}. For
1-hop queries, whether executed using oblivious relational operators or custom query 
processing algorithms, the output or intermediary sizes are bounded by $|e|$, which
is considered non-sensitive. However, expanding the system to support
multi-hop queries introduces the challenge of hiding intermediary result sizes.
Naively addressing this by applying worst-case padding — pad an intermediate result to its worse possible size — significantly degrades system performance. 
For example,
for a 2-hop query spanning edges $e_1$ and $e_2$, worst-case padded intermediary table
will be of size $|e_1| \times |e_2|$.


\begin{tcolorbox}[colback=gray!7!white,coltitle=black,colframe=orange!50!white,title=\centering \textbf{Research question 3}]
\centering
\emph{Can we develop more \textbf{efficient} mechanisms to hide or obfuscate intermediary result sizes without resorting to worst-case padding?}
\end{tcolorbox}

\noindent\textit{\textbf{Towards a solution}}: Computing the size of any table is 
equivalent to computing the COUNT aggregate function on a table. Differential privacy (DP) \cite{DworkRoth2014} is an ideal technique for hiding the exact results of an aggregate query. Shrinkwrap \cite{Bater2018} was
the seminal work that applied DP to hide or obfuscate intermediary result sizes of
a SQL query consisting of multiple operators. However, it considered a privacy budget
$\epsilon$ allocated to each query executed by the system. Thus, the \textit{global 
privacy budget} increases significantly per executed query, potentially revealing more
than anticipated amount of sensitive information. Alternatively, solutions such as 
Turbo \cite{Kostopoulou2023} and cacheDP \cite{Mazmudar2022} cache the DP-derived noise for queries and reuse them when queries
are repeated to reduce the consumption of the global privacy budget for each query. 
However, the proposed techniques only apply for \textit{linear queries} where queries 
span a single table. As discussed so far, queries in graph databases cover many nodes
and edges, and the process of joining them amplifies query sensitivity, thereby increasing 
the required noise budget. Hence, this challenge presents new research opportunities in 
caching DP-noise for complex queries that span multiple entities or tables.

\vspace{-0.7em}
\subsection{Cyclic queries}
\label{sub:cyclic}

A common query pattern observed in property graph databases are cyclic
queries of the form $n1 - e1 - n2 \- \cdots e_i - n1$.

\noindent\textbf{\textit{Challenges with cyclic queries}}. While cyclic queries can be 
executed with a 
series of binary joins, such an approach often generates 
exponentially large intermediate results~\cite{ngo2014skew}, significantly 
degrading performance in oblivious settings due to their inherent computational overhead.
For instance, consider the query: "users who are friends and have purchased the same product." This forms a cyclic pattern: (u1:User) - [:FRIEND] → (u2:User) - [:PURCHASE] → (:Product) ← [:PURCHASE] - (u1:User). A 
binary join approach might first compute all User-Friend pairs (potentially $O(U^2)$ for $U$ users), then join with Purchases to find shared products. For example, with 1,000 users, each having 100 friends and 10 purchases, the User-Friend join alone produces 100,000 pairs, which, when joined with Purchases, results in 1,000,000 intermediate tuples—most of which are discarded later.

\begin{tcolorbox}[colback=gray!7!white,coltitle=black,colframe=orange!50!white,title=\centering \textbf{Research question 4}]
\centering
\emph{Can we device oblivious join algorithms for complex cyclic joins that are
asymptotically \textbf{faster} than a sequence of binary joins?}
\end{tcolorbox}

\noindent\textit{\textbf{Towards a solution}}:
Worst case optimal joins (WCOJ)~\cite{ngo2018worst} handle cyclic queries more efficiently by evaluating multi-way intersections column-at-a-time rather than table-at-a-time, achieving the theoretical worst case output size defined by the AGM bound~\cite{atserias2008size} (e.g., $n^{3/2}$ for a cycle of 3). In the example above, processing the query using WCOJ 
simultaneously tracks the user-friend relation and the intersecting purchases, which prunes non-matching paths early. While a few oblivious acyclic
multi-way joins have been proposed in the past~\cite{arasu2013oblivious, chu2021differentially}, Hu and Wu~\cite{hu2025optimal} recently
proposed the first oblivious WCOJ construction. However, their construction pads the output to a query-dependent worst-case
size (i.e., $O(n^f)$ where $f$ indicates fractional edge cover), which we believe
can be further reduced. Additionally, all the above discussed papers~\cite{arasu2013oblivious, chu2021differentially,hu2025optimal} are theoretical in nature with no implementation to show empirical results. Hence, developing practical oblivious
WCOJs for a secure property graph database remains an open area of research. 
% Property graphs benefit from WCOJ's ability to handle cyclic many-to-many relationships, it is common in recommendation systems (e.g. friends who liked the same items) or fraud detection (e.g. cycles of transactions).  

